% Part: computability
% Chapter: computability-theory
% Section: russells-paradox

\documentclass[../../../include/open-logic-section]{subfiles}

\begin{document}

\olfileid{cmp}{thy}{rus}
\olsection{مقارنة بمفارقة راسل}

ويتضح شأن هذه الحجج بموازنتها بمفارقة راسل، لنرى ما يأتلف منها وما
يفترق:
\begin{enumerate}
\item مفارقة راسل: لتكن $S=\Setabs{x}{x\notin x}$. عندئذ
$S\in S$ إذا وفقط إذا كان $S\notin S$، وهذا تناقض.

  \emph{النتيجة:} يمتنع وجود المجموعة~$S$ على هذه الصفة؛ فإثبات «مجموعة جميع
  المجموعات» يناقض سائر بديهيات نظرية المجموعات.

\item تعديل لمفارقة راسل: لتكن $F$ «الدالة» من مجموعة جميع الدوال إلى
$\{0,1\}$، المعرفة بـ
  \[
  F(f) =
  \begin{cases}
    1 & \text{إذا كانت $f$ في مجال $f$ وكان $f(f) = 0$} \\
    0 & \text{في غير ذلك}
  \end{cases}
  \]
  وبمثل الحجة الأولى يلزم أن $F(F)=0$ إذا وفقط إذا كان $F(F)=1$، وهذا تناقض.

  \emph{النتيجة:} لا تكون $F$ دالة؛ إذ «مجموعة جميع الدوال» أعظم من أن
  يصح اتخاذها مجالًا لدالة.

\item الحجة القطرية: لتكن $f_0$ و$f_1$ و… تعداد الدوال الجزئية
  القابلة للحساب، ولتكن $G\colon\Nat\to\{0,1\}$ معرفة بـ
  \[
  G(x) =
  \begin{cases}
    1 & \text{إذا كان $f_x(x)\downarrow = 0$} \\
    0 & \text{في غير ذلك}
  \end{cases}
  \]
  فلو كانت $G$ قابلة للحساب، لكانت $f_k$ لعدد ما $k$، وحينئذ
  يكون $G(k)=1$ إذا وفقط إذا كان $G(k)=0$، وهذا تناقض.

  \emph{النتيجة:} تمتنع قابلية $G$ للحساب، ولا يمتنع أن تكون $G$ دالة في
  نظرية المجموعات. فليس هذا تناقضًا في البديهيات، وإنما هو تمييز لما يقبل الحساب مما لا يقبله.
\end{enumerate}

وقد يلتبس تداخل هذه الأسماء: الدوال الجزئية، والدوال القابلة للحساب،
والدوال الجزئية القابلة للحساب. فلنميزها. الدوال
الجزئية من $\Nat$ إلى $\Nat$ تؤلف مجموعة واسعة؛ منها الكلي،
ومنها القابل للحساب، ومنها ما يجمع الأمرين، ومنها ما يفقدهما
معًا. واصطلاحنا، حين نطلق «دالة» بلا قيد، أن نريد الكلية.
فتكون الأصناف كما يأتي:
\begin{enumerate}
\item دوال قابلة للحساب؛
\item دوال جزئية قابلة للحساب وليست كلية؛
\item دوال غير قابلة للحساب؛
\item دوال جزئية ليست كلية ولا قابلة للحساب.
\end{enumerate}
ولتستحضر هذه الأصناف، ارسم مربعًا يمثل جميع الدوال الجزئية من
$\Nat$ إلى $\Nat$، واجعل فيه منطقتين متداخلتين: إحداهما للدوال الكلية،
والأخرى للجزئية القابلة للحساب. ثم تمرّن بأن
تصف دالة في كل منطقة ينقسم إليها الرسم.

\end{document}
